% section: setup
\section{Experimental setup}
\label{sec:setup}

Every number in this section was verified against a primary source or measured directly; the
provenance of each is recorded in \texttt{verification/phase0.md} in the accompanying repository.
Where a figure differs from the one in our original design document, the measured value is given and
the discrepancy is stated explicitly rather than silently corrected.

\subsection{Model}
\label{sec:setup:model}

The backbone is \texttt{Qwen/Qwen3-VL-2B-Instruct} (Apache-2.0, ungated). Two properties make it
suited to this task rather than merely small enough to train on free hardware. First, its published
ChartQA test score of \result{79.1} gives a known, matched starting point
(Qwen3-VL technical report, Table~4: \result{86.6} Thinking / \result{79.1} Instruct). Second, and
more importantly for the grounding half of this work, Qwen3-VL emits bounding boxes in a
\emph{normalised $0$--$1000$ coordinate system}, changed from the absolute pixel coordinates of
Qwen2.5-VL. RefChartQA's evaluator quantises to $0$--$999$, so almost no conversion is required and
correspondingly few places exist for a coordinate bug to hide.

\paragraph{Visual tokenisation.}
Qwen3-VL uses \texttt{patch\_size} $=16$ and \texttt{spatial\_merge\_size} $=2$, so one visual token
covers a $32\times32$ pixel region. This corrects a figure we had inherited: the widely quoted
$28\times28$ token is the Qwen2-VL/Qwen2.5-VL geometry (patch $14\times$ merge $2$). The correction
is not cosmetic. A \emph{larger} minimum visual unit means \emph{more} grounding targets fall below
it, which strengthens rather than weakens the central premise of this work
(Section~\ref{sec:setup:subtoken}).

We derive the factor and the pixel bounds from the loaded processor at runtime rather than
hard-coding them, and cross-check our \texttt{smart\_resize} port against
\texttt{transformers.models.qwen2\_vl.image\_processing\_qwen2\_vl.smart\_resize} across ten image
shapes at both factors.

\subsection{The sub-token problem}
\label{sec:setup:subtoken}

A model cannot reliably localise a region smaller than its own smallest visual unit. Chart evidence
regions frequently are smaller. Measured over \result{800} RefChartQA validation images
(\result{1{,}045} boxes; the test split is sealed until Section~\ref{sec:results}):

\begin{center}
\begin{tabular}{lrr}
\toprule
Input budget & Median visual tokens & Targets sub-token \\
\midrule
$448^2$        & \result{176} & \result{60.5\%} \\
$512^2$        & \result{247} & \result{53.2\%} \\
$640^2$        & \result{368} & \result{43.1\%} \\
$768^2$ (native) & \result{425} & \result{41.3\%} \\
\bottomrule
\end{tabular}
\end{center}

``Sub-token'' here means the target is narrower than one visual token on at least one axis. Two
observations follow. At the $512$-pixel budget originally planned, \emph{more than half} of all
grounding targets are unresolvable in principle. And because these charts are only about $800$
pixels wide, any budget at or above $768^2$ is already native: moving from $512^2$ to native costs
$1.7\times$ the tokens, not the $4\times$ a naive analysis predicts, and recovers roughly twelve
percentage points of targets. Input resolution is therefore treated as a measured variable rather
than a constant (Section~\ref{sec:setup:compute}).

\subsection{Datasets}
\label{sec:setup:data}

\begin{center}
\begin{tabular}{llrl}
\toprule
Dataset & Split & Rows & Licence \\
\midrule
ChartQA     & train / val / test & \result{28{,}299} / \result{1{,}920} / \result{2{,}500} & GPL-3.0 \\
RefChartQA  & train / val / test & \result{55{,}789} / \result{6{,}223} / \result{11{,}690} & AGPL-3.0 \\
ChartQAPro  & test only          & \result{1{,}948} & MIT \\
\bottomrule
\end{tabular}
\end{center}

RefChartQA's test split partitions by question source into \result{500} human, \result{1{,}032}
machine and \result{10{,}158} program-of-thought items. This matters more than it appears: the
official evaluator produces \emph{three separate scores} and no aggregate, and the published
reference figure of \result{32.83} AP@0.5 that this work targets is the \emph{human} subset alone.
Comparing a whole-test-set average against it would not be a comparison. All grounding results in
this report are therefore given as a per-subset triple, and the comparison against the published
figure is stated as being on RefChartQA-H only.

The consequence is a hard limit on resolution: with $n=500$, a 95\% bootstrap interval on a score in
the low thirties spans roughly $\pm4$ points, so differences smaller than that are not
distinguishable from sampling noise. We state this before reporting any result.

ChartQA's gold data tables, required for plan mining, are \emph{not} present in the parquet
distribution; they live in a separate \result{875{,}370{,}872}-byte archive in the same repository.

\subsection{Metrics, and the implementations that define them}
\label{sec:setup:metrics}

A benchmark is a procedure, not a dataset: a score is a property of the model \emph{together with}
the prompt, the decoding settings, the output format, the parsing code, the metric implementation
and the data version. We therefore fix the implementations explicitly.

\paragraph{Answer accuracy.}
The ChartQA repository contains no answer evaluator. The implementation every published ChartQA
number was produced with is \texttt{relaxed\_correctness} from pix2struct, which RefChartQA vendors
verbatim; we verified the two agree on all \result{216} pairs of a $12\times18$ target/prediction
grid. One function therefore scores both of our protocols.

That function has a quirk worth naming, because it is canonical rather than incidental: the guard
reads \texttt{if prediction\_float is not None and target\_float}, a \emph{truthiness} test, so a
gold answer parsing to $0.0$ is falsy and the comparison silently falls through to case-insensitive
string equality. A prediction of \texttt{"0.0"} against a gold \texttt{"0"} is therefore scored
incorrect. Every published ChartQA number carries this behaviour, so we adopt it unchanged rather
than ``fixing'' it into incomparability, and normalise our answers to bare \texttt{"0"}.

\paragraph{Grounding.}
AP@0.5 and P@F1 come from RefChartQA's released evaluator, vendored byte-identically and pinned by
hash. Characterising it by execution rather than by reading revealed two properties that materially
constrain how a system should emit boxes:

\begin{enumerate}
  \item Every predicted box is assigned $\mathrm{score}=1.0$, so there is no confidence ranking and
        AP degenerates to $1/(\text{rank of the first correct box})$. A single false positive placed
        ahead of a true one halves AP and sends P@F1 to zero.
  \item \texttt{compute\_AP\_50} pools every prediction across the dataset into a single
        precision--recall curve. Extra boxes are consequently \emph{free} on a single image and
        ruinous in aggregate: over twenty images, appending three false positives to an otherwise
        perfect prediction set moves AP from \result{1.0000} to \result{0.3243}.
\end{enumerate}

The system therefore emits few boxes, ordered best-first, and the evidence list is filtered before
scoring under a rule frozen on validation data.

A further trap: the evaluator accepts a box only when all four coordinates lie in $[0, 999]$, and
\emph{silently discards the entire box otherwise}. Qwen3-VL's native range is $[0,1000]$ inclusive,
and boxes touching a chart's right or bottom edge are common, so all emitted coordinates are clamped.

\subsection{Compute}
\label{sec:setup:compute}

All training and evaluation runs on free-tier hardware: a single NVIDIA T4 (about $15$\,GiB) on
Kaggle, with Colab as a secondary host. No paid compute is used. Runs are resumable and push
checkpoints to a private repository on every save, because free sessions are terminated without
warning.

\TODO{fill from smoke\_results.json once Phase 2 completes: peak reserved memory and seconds per
optimizer step for each backend and input budget, and the resulting backbone and resolution choice}

% tab_compute — populated by: cdt-report --what compute
% Source of record: verification/measured_facts.json and RUNS.md.
\begin{table}[t]
\centering
\caption{Compute characterisation on free-tier hardware. All figures are for a
\emph{single} NVIDIA T4; a run whose model is sharded across two cards is rejected rather than
reported, because the budget these gates encode is a one-card budget
(Section~\ref{sec:setup:compute}).}
\label{tab:compute}
\begin{tabular}{lrrr}
\toprule
Configuration & Peak reserved (GiB) & s / optimizer step & Projected 3{,}000 steps (h) \\
\midrule
\TODO{512px, single card} & & & \\
\TODO{native, single card} & & & \\
\midrule
Gate & $\leq 13.5$ & --- & $\leq 10$ \\
\bottomrule
\end{tabular}
\end{table}

